\documentclass{jarticle}
\usepackage[dvipdfmx]{graphicx}
\topmargin       -0.4mm
\oddsidemargin   -0.4mm
\evensidemargin  -0.4mm
\headheight       0.0mm
\headsep          0.0mm
\footskip        10.0mm
%
\textheight     247.0mm
\textwidth      170.0mm
%
\setlength{\parindent}{1zw}
%
\usepackage{graphicx}
\usepackage{multirow}
\usepackage{wrapfig}
\makeatletter
\def \fgcaption{\def \@captype{figure} \caption}
\def \tbcaption{\def \@captype{table} \caption}
\makeatother
\def \figurename{Fig.}
\def \tablename{Table}
%
\def\topfraction{.95}
\def\bottomfraction{.95}
\def\textfraction{.05}
\def\floatpagefraction{.8}
\def\dbltopfraction{.95}
\def\dbltextfraction{.05}
\def\dblfloatpagefraction{.8}
%
\pagestyle{empty}
%
\begin{document}
%
\begin{center}
 \begin{tabular}{|c|c|c|} \hline
  \parbox{50mm}{\vspace{1mm} \centering{\Large\bf 202*/**/**}} & 
  \parbox{55mm}{\vspace{1mm} \centering{\Large\bf B4　姓名　○○}} & 
  \parbox{50mm}{\vspace{1mm} \centering{\Large\bf No.~1}}
  \\ [2mm] \hline
%
  \multicolumn{3}{|c|}{%
   \parbox{155mm}{\centering\Large\bf
    \vspace{2mm}
    ＊＊＊の研究（研究題目）% ここに研究題目を書く
   }
  }
  \\ [3mm] \hline\hline
%
  \multicolumn{3}{|c|}{%
   \parbox{155mm}{\large\vspace{2mm}
%
    {\Large\bf 今月の目標}
%
    \begin{itemize}
     \item ＊＊に関する計測実験とデータ解析を完了
    \end{itemize}
%
    % 進捗の要点(箇条書)
    % 前回の要点と同じものを挙げない
    {\Large\bf 進捗の要点}
%
    \begin{itemize}
     \item 今回進捗した要点を列挙する，前回の要点と同じものを挙げない
    \end{itemize}
%
    % 今後の実施事項(箇条書)
    % 前回の事項と同じものを挙げない
    {\Large\bf 今後の実施事項}
%
    \begin{itemize}
     \item 次回までに実施する事項を列挙する，前回の事項と同じものを挙げない
    \end{itemize}
%
   }
  }
  \\ [3mm] \hline
 \end{tabular}
%
\end{center}
%
% ここから下は進捗内容を具体的に書くこと．
% 例えば，実験を行ったのなら，その方法と結果を載せる．
% また，アルゴリズムを考えた場合は，そのアルゴリズムの詳しい説明を書く．
%
%\par

この報告書では，月初めの回で月末までの目標を説明する．\\以降の回では前回からの進捗を報告する．次の点について明記する．
\begin{enumerate}
\item 何をどのように進めたのか？
\item それによって，どのような結果が得られたのか？
\item その結果がどのような意味を持つのか？
\item 次回報告までの今後の検討事項
\end{enumerate}
\par
グラフ等は例えばFig.~\ref{fig:graph1}に示すように図を使い説明する．参考にした本，論文があれば引用する\cite{ito1991}．
%
\begin{figure}[htbp]
  \begin{center}
   \includegraphics [scale=0.5]{./plot.png}	% jpgでも可
  \end{center}
\vspace{-7mm}
\caption{Difference of reaction force.}
\label{fig:graph1}
\end{figure}
%

注）このひな形は，TeX liveをインストールした環境において，pdf化が可能である．
%

\begin{thebibliography}{9}
\bibitem{ito1991}
伊藤和人, ''\LaTeX トータルガイド'', 秀和システムトレーディング, 1991.
\end{thebibliography}
%
%
\end{document}